\documentclass[11pt]{article}
\usepackage[utf8]{inputenc}
\usepackage[T2A]{fontenc}
\usepackage[russian]{babel}
\usepackage{amsmath,amssymb,amsthm}
\usepackage[margin=2.5cm]{geometry}
\usepackage{hyperref}
\usepackage{booktabs}

\newtheorem{definition}{Определение}
\newtheorem{property}{Свойство}

\title{Формализация MPC-слоя: PRSS, threshold-HE, compartmented secret sharing}
\author{}
\date{11 сентября 2026}

\begin{document}
\maketitle

\section{Постановка}

Три пункта из рабочих заметок формализуются по отдельности: (1) протокол суммирования
$N$ участников, в котором сервер \emph{не} входит в множество участников и не
получает ключевого материала; (2) гомоморфное шифрование как альтернатива/дополнение
к маскированию; (3) поправка модели угроз на то, что участники одного домохозяйства
не являются независимыми доверительными единицами (\guillemotleft знают двое —
знает семья\guillemotright).

Обозначения: $P = \{1,\dots,N\}$ — множество роботов-участников, $S$ — сервер,
$S \notin P$. Каждый участник $i$ хранит приватное значение $x_i$ (например,
non-conformity score или дельту весов). Цель: $S$ получает $\sum_i x_i$ и ничего
сверх этого.

\section{Механизм 1: PRSS — суммирование без сервера как участника}

\begin{definition}[Pseudorandom Secret Sharing]
Пусть $\Gamma$ — $(t,N)$-пороговая структура доступа над $P$. PRSS-схема
\cite{prss2005} позволяет каждому участнику $i$, зная только заранее (один раз, до
начала обучения) распределённые PRF-ключи, локально и без коммуникации в каждом
раунде вычислять свою долю свежего Shamir-разделения нуля $z_i(\mathrm{sid})$ для
произвольного идентификатора раунда $\mathrm{sid}$, такую что для любого
квалифицированного (по $\Gamma$) подмножества $Q\subseteq P$ значения
$\{z_i(\mathrm{sid})\}_{i\in Q}$ лежат на многочлене степени $t-1$ с $p(0)=0$.
\end{definition}

Протокол на раунде $\mathrm{sid}$:
\begin{enumerate}
  \item (Однократная настройка, до обучения.) Для каждого максимального
        неквалифицированного множества $B\subset P$, $|B|=t-1$, все участники
        $P\setminus B$ получают общий PRF-ключ $k_B$. Участник $i$ хранит
        $\{k_B : B \text{ макс.\ неквал.}, i\notin B\}$.
  \item Участник $i$ вычисляет маску
        $z_i(\mathrm{sid}) = \sum_{B:\,i\notin B} c_{i,B}\cdot \mathrm{PRF}(k_B,\mathrm{sid})$,
        где коэффициенты $c_{i,B}$ фиксированы на этапе настройки и не зависят от
        $\mathrm{sid}$.
  \item Участник отправляет серверу $y_i = x_i + z_i(\mathrm{sid})$.
  \item Сервер вычисляет $\sum_{i\in P} y_i$.
\end{enumerate}

\begin{property}[Корректность]
При полном участии всех $N$ сторон $\sum_{i} z_i(\mathrm{sid}) = 0$ по построению
PRSS-схемы \cite{prss2005}, следовательно $\sum_i y_i = \sum_i x_i$.
\end{property}

\textbf{Почему сервер формально исключён из участников:} структура доступа $\Gamma$
и PRF-ключи $\{k_B\}$ определены только над $P$; $S$ не получает ни одного ключа ни
на каком этапе. Даже при полной компрометации $S$ единственное, что можно извлечь —
это переданные $y_i$ и их сумму; восстановить отдельные $x_i$ невозможно без доступа
к ключам, которых у $S$ никогда не было.

\textbf{Устойчивость к отвалившимся участникам.} Базовая PRSS-конструкция
гарантирует sum-to-zero только для полного множества $P$; при частичном участии
(часть роботов офлайн) нужен дополнительный слой — восстановление вклада отпавших
через Shamir-разделение их собственных PRF-семян среди оставшихся (как в
\cite{bonawitzref}), либо специализированная straggler-resilient PRSS-конструкция
\cite{prssstraggler2021}, прямо решающая эту проблему для нестабильной связи —
актуально для домашних роботов на Wi-Fi.

\section{Механизм 2: гомоморфное шифрование (threshold-HE)}

Альтернатива маскированию — аддитивно гомоморфное шифрование (Paillier
\cite{paillier1999}) с ключом расшифровки, разделённым по Shamir среди участников.

\begin{enumerate}
  \item (Настройка.) Публичный ключ $\mathrm{pk}$ известен всем и серверу; секретный
        ключ $\mathrm{sk}$ разделён на доли $\mathrm{sk}_1,\dots,\mathrm{sk}_N$ по
        схеме Shamir с порогом $t$. Сервер \emph{не} получает ни $\mathrm{sk}$, ни
        ни одной доли.
  \item Участник $i$ шифрует: $c_i = \mathrm{Enc}(\mathrm{pk}, x_i)$.
  \item Сервер (или любой агрегатор) вычисляет
        $c = \prod_{i=1}^N c_i \bmod n^2 = \mathrm{Enc}\big(\mathrm{pk}, \textstyle\sum_i x_i\big)$
        по аддитивному гомоморфизму Paillier — без расшифровки отдельных $c_i$.
  \item Расшифровать $c$ можно только коллективно: $\ge t$ участников вычисляют
        частичные расшифровки $d_i = \mathrm{PartialDec}(\mathrm{sk}_i, c)$ и
        комбинируют их \cite{thresholdpaillier}.
\end{enumerate}

\begin{center}
\begin{tabular}{@{}p{3.4cm}p{5.6cm}p{5.6cm}@{}}
\toprule
 & PRSS-маскирование & Threshold-HE (Paillier) \\
\midrule
Основа безопасности & информационно-теоретическая (истинная случайность PRF) &
  вычислительная (decisional composite residuosity) \\
Стоимость на значение & сложение по модулю (дёшево, годится для слабого бортового
  компьюта) & модулярное возведение в степень (заметно дороже) \\
Когда нужны $\ge t$ живых участников & на этапе маскирования (каждый раунд) & только
  на этапе финальной расшифровки \\
Устойчивость к отвалившимся & нужен отдельный dropout-recovery слой & агрегация
  ciphertext'ов возможна асинхронно, decrypt откладывается до сбора $t$ участников \\
\bottomrule
\end{tabular}
\end{center}

Практическое следствие для флота домашних роботов с нерегулярным расписанием
работы: threshold-HE позволяет накапливать зашифрованные вклады асинхронно (робот
шифрует и отправляет, когда он online) и отложить сам момент расшифровки до
удобного окна, когда $\ge t$ роботов одновременно доступны — тогда как
PRSS/маскирование по Bonawitz-схеме требует синхронизации на каждом раунде
маскирования.

\subsection*{PRSS и Paillier не взаимоисключающие — их можно совместить}

Вопрос, можно ли объединить маскирование (sum-to-zero) и гомоморфное шифрование,
имеет прямой ответ: да, и это не самодеятельность — устоявшееся направление
(\emph{LERNA} \cite{lerna2023}, ASIACRYPT 2023; прикладные FL-варианты
\emph{DMSA-FL} \cite{dmsafl2026} и \emph{V-MHESA} \cite{vmhesa2025}).

\textbf{Конструкция.} Вместо шифрования сырого значения участник шифрует его вместе
с PRSS-маской:
\[
  c_i = \mathrm{Enc}\big(\mathrm{pk},\; x_i + z_i(\mathrm{sid})\big), \qquad
  \sum_i z_i(\mathrm{sid}) = 0.
\]
По аддитивному гомоморфизму:
\[
  \prod_i c_i = \mathrm{Enc}\Big(\mathrm{pk},\; \sum_i x_i + \sum_i z_i(\mathrm{sid})\Big)
  = \mathrm{Enc}\Big(\mathrm{pk},\; \sum_i x_i\Big),
\]
маска сокращается тем же способом, что и в чистом PRSS-механизме (раздел~2), а
корректность агрегата не зависит от того, что значения дополнительно зашифрованы.

\textbf{Зачем это нужно, если Paillier и так не раскрывает $x_i$.} Сам по себе
threshold-Paillier ничего не говорит о том, \emph{что именно} разрешено
расшифровывать коалиции из $t$ участников — протокол защищает секретный ключ, но не
запрещает применить его к отдельному $c_i$ вместо произведения $\prod c_i$. Если
$t$ участников (или сговорившийся сервер с $t$ участниками) решат расшифровать не
агрегат, а конкретный $c_i$ одного робота, threshold-HE сам по себе это не
остановит — по протоколу это валидная операция. PRSS-маска под шифром закрывает
именно эту дыру: расшифровка одного $c_i$ отдельно от остальных даёт
$x_i+z_i(\mathrm{sid})$, а не $x_i$, и $z_i$ защищена отдельной структурой PRF-ключей
(раздел~2), не входящей в threshold-Paillier ключ. Тем самым комбинация закрывает
и внешнюю угрозу (наблюдатель шифроканала без ключа — защищает HE), и внутреннюю
(коалиция с ключом, пытающаяся расшифровать не то, что положено — защищает маска).

\textbf{Статус.} В отличие от домохозяйственного уточнения ниже, это уже не
предлагаемая, а существующая в литературе конструкция — прямо переносится на MPC-слой
проекта без дополнительных исследований, только инженерная интеграция.

\section{Механизм 3: \guillemotleft знают двое — знает семья\guillemotright{} —
compartmented secret sharing}

\subsection*{Проблема наивного $(t,N)$-Shamir}

Плоская пороговая схема считает всех $N$ участников независимыми доверительными
единицами: любые $t$ долей восстанавливают секрет, любые $t-1$ — нет, вне
зависимости от того, \emph{кому принадлежат} эти доли. Но если два робота
принадлежат одному домохозяйству, они физически не независимы: общий Wi-Fi, общий
владелец аккаунта, физический доступ к обоим устройствам одним и тем же человеком.
Компрометация одного робота дома практически означает доступность и второго —
\guillemotleft знают двое (робота одной семьи) — знает семья\guillemotright,
не два независимых доверительных агента.

\subsection*{Формализация через compartmented access structure}

\begin{definition}[Compartmented threshold structure, Simmons 1990 / Brickell 1990
  \cite{compartmented1990}]
Множество участников $P$ разбито на непересекающиеся группы (compartments)
$H_1,\dots,H_m$, $P = \bigsqcup_j H_j$ (в нашем случае $H_j$ — домохозяйство).
Задаются пороги $k_1,\dots,k_m$ по компартментам и, при необходимости, общий
верхнеуровневый порог. Секрет восстановим тогда и только тогда, когда для каждой
участвующей группы $H_j$ набралось $\ge k_j$ членов.
\end{definition}

Для анализа приватности (не восстановления, а \emph{угрозы} восстановления)
релевантна двойственная величина — \textbf{домохозяйственный след} колеблющегося
множества $S\subseteq P$:
\[
  h(S) = \bigl|\{\, j : H_j \cap S \ne \varnothing \,\}\bigr|.
\]

\begin{property}[Поправка порога]
Наивный порог $|S|\ge t$ переоценивает реальную независимость коалиции, если
$H_j$ содержат более одного участника. Корректное условие приватности: секрет
защищён, пока $h(S) < \tau$ для некоторого порога по числу \emph{независимых}
домохозяйств $\tau$, а не пока $|S| < t$ по сырому числу роботов.
\end{property}

Иными словами, при плоском $(t,N)$-Shamir с $t=5$ атакующий, скомпрометировавший
всего 2 домохозяйства по 3 робота в каждом ($|S|=5$, но $h(S)=2$), формально
достигает порога $t$, хотя реально контролирует только 2 независимых источника
доверия, а не 5. Правильная конструкция — присвоить каждому домохозяйству единый
слот верхнеуровневого секрета (внутренний порог $k_j$ внутри дома, почти всегда
тривиальный при 1--3 роботах на дом, то есть $k_j=1$), и требовать порог $\tau$
именно по числу участвующих домохозяйств, как в двухуровневой конструкции
Brickell \cite{compartmented1990}.

\textbf{Статус.} Прямого применения compartmented secret sharing к флотам бытовых
роботов в проверенной литературе не найдено — это предлагаемое уточнение модели
угроз проекта, не established результат; встраивается в byzantine/групповой гэп,
уже отмеченный в матчасти по Shamir/MPC.

\section{Механизм 4: разделение между несколькими серверами}

Все механизмы выше предполагали \emph{один} сервер $S$, не входящий в множество
участников. Альтернатива — разделить сам сервер на $m$ независимых серверов
$S_1,\dots,S_m$, ни один из которых по отдельности не видит агрегат, если они не
сговариваются друг с другом. Два эталонных образца:

\begin{itemize}
  \item \textbf{Prio} \cite{prio2017} (Corrigan-Gibbs, Boneh, NSDI 2017) — $m=2$
        не сговаривающихся сервера. Клиент режет значение на 2 аддитивные доли, по
        одной каждому серверу; пока хотя бы один сервер честный, ничего не
        раскрывается до момента, когда серверы совместно соглашаются опубликовать
        агрегат.
  \item \textbf{ABY3} \cite{aby3_2018} — эталон для $m=3$ с честным большинством
        (не более 1 из 3 нечестен), replicated $2$-из-$3$ secret sharing: каждому
        серверу достаётся 2 из 3 долей каждого значения, что даёт устойчивость к
        выходу одного сервера из строя.
\end{itemize}

\subsection*{Конструкция для суммы (наш случай)}

Для линейной функции (суммы) разделение между $m=3$ серверами не требует вообще
никакой криптографии на стороне робота — только аддитивное разрезание числа.
Робот $i$ представляет своё значение как
\[
  x_i = x_i^{(1)} + x_i^{(2)} + x_i^{(3)} \pmod p,
\]
случайно выбирая первые два слагаемых и отправляя $x_i^{(k)}$ серверу $S_k$. Каждый
сервер \textbf{локально}, без единого раунда взаимодействия с другими серверами,
складывает свою колонку по всем $N$ роботам:
\[
  S^{(k)} = \sum_{i=1}^N x_i^{(k)}, \qquad k=1,2,3.
\]
Итоговый агрегат $\sum_i x_i = S^{(1)}+S^{(2)}+S^{(3)}$ собирается только на
последнем шаге, объединением трёх частичных сумм; до этого момента ни один сервер
по отдельности не располагает ничем, кроме своей доли.

\subsection*{Сравнение с single-server механизмами (1--3)}

\begin{center}
\begin{tabular}{@{}p{3.6cm}p{5.6cm}p{5.6cm}@{}}
\toprule
 & Один сервер (PRSS/HE, разделы 1--2) & Несколько серверов (Prio/ABY3) \\
\midrule
Основа доверия & криптографическая (PRF-ключи / decisional composite
  residuosity) & организационная (\guillemotleft серверы не сговариваются\guillemotright) \\
Стоимость на робота & сложение (PRSS) или модулярная экспонента (HE) & только
  аддитивное разрезание числа — дешевле обоих \\
Единая точка отказа/атаки & да, один сервер — цель атаки & нет, атакующему нужно
  скомпрометировать несколько независимых операторов сразу \\
\bottomrule
\end{tabular}
\end{center}

\textbf{Ключевая оговорка.} Гарантия здесь не математическая, а организационная:
если все $m$ серверов принадлежат одному вендору (пусть и в разных дата-центрах),
допущение о несговоре фиктивно. Требуется реально независимый второй/третий
оператор (партнёр, регулятор, независимый аудитор) — что прямо стыкуется с
нарративом проекта про 152-ФЗ/GDPR и распределение данных по юрисдикциям/
организациям, а не просто добавляет крипто-сложности.

\section{Направление улучшения: verifiable weighted aggregation + group testing}

По запросу \guillemotleft хорошенько изучить математику процесса и улучшить
её\guillemotright{} — через arXiv MCP найдены три работы 2025--2026 годов, которые
вместе дают конкретный, более сильный план вместо схемы compartmented 1990 года из
раздела выше.

\subsection*{Verifiable Weighted Secret Sharing вместо compartmented}

\emph{Verifiable Weighted Secret Sharing} \cite{vwss2025} (Shehata, Han,
Thyagarajan, 2025) обобщает ровно ту проблему, которую решает раздел выше, но
современным и эффективным инструментом: вместо равных долей каждому участнику
присваивается \textbf{вес} $\omega_i$, и безопасность гарантируется, пока суммарный
вес скомпрометированных участников не превышает порог $t$ — то есть
$\sum_{i\in S} \omega_i < t \Rightarrow$ секрет защищён. Это формально сильнее и
гибче предложенной выше двухуровневой compartmented-конструкции: вместо
\guillemotleft каждому домохозяйству — единый слот\guillemotright{} достаточно
задать $\omega_i = 1/|H_{j(i)}|$ для робота $i$ из дома $H_{j(i)}$ — тогда сумма
весов всех роботов одного дома равна ровно 1 (\guillemotleft знают двое — знает
семья\guillemotright{} становится буквальным следствием весовой формулы, а не
отдельной надстройкой). Схема к тому же \textbf{верифицируема} — защищена не только
от honest-but-curious, но и от злонамеренного дилера — и на практике даёт
100$\times$ улучшение по коммуникационной сложности относительно наивной
взвешенной VSS и 20$\times$ относительно невзвешенной VSS (числа из оценки на
Ethereum \cite{vwss2025}, но конструкция общего назначения, применима и вне
блокчейна).

\subsection*{Group testing для byzantine-обнаружения поверх той же группировки}

\emph{FedGT} \cite{fedgt2023} (Xhemrishi et al., 2023) решает соседнюю, но не ту же
задачу: обнаружение вредоносных (byzantine) клиентов \emph{внутри} secure
aggregation, через пересекающиеся группы клиентов и декодирование по методу group
testing. Сервер узнаёт агрегат по каждой группе (не по каждому клиенту и не по
всему флоту сразу), что позволяет вычислить, какие группы содержат вредоносного
участника, и исключить их поставщиков из обучения — не жертвуя приватностью
отдельных участников.

\textbf{Ключевое наблюдение:} группы в FedGT можно естественно выбрать
\emph{совпадающими с домохозяйствами} $H_1,\dots,H_m$ (плюс дополнительные
пересекающиеся группы для самого group-testing декодирования). Тогда одна и та же
группировка по домам одновременно: (a) определяет веса в verifiable weighted
aggregation из предыдущего пункта, и (b) служит структурой для byzantine-детекции
FedGT — закрывая byzantine-гэп, отмеченный в разделе про Shamir/MPC, тем же
организационным примитивом, а не отдельным механизмом.

\subsection*{Граница эффективности по размеру доли}

\emph{Optimal Computational Secret Sharing} \cite{optcss2025} (Aureliano, Cohen,
D'Oliveira, 2025) даёт информационно-теоретическую нижнюю границу и достигающую её
конструкцию для размера доли при вычислительной (а не только
информационно-теоретической) секретности: $\frac{|S|+|K|}{t}$ вместо $|S|$ для
классического Shamir — где $|K|$ — размер криптографического ключа. Это верхний
ориентир эффективности, к которому стоит стремиться при выборе конкретной
weighted/compartmented конструкции для флота: доля на робота должна
асимптотически приближаться к этой границе, а не быть полным размером секрета,
что критично при ограниченном бортовом хранилище.

\subsection*{Предлагаемый план}

\begin{enumerate}
  \item Заменить плоский $(t,N)$-Shamir проекта на CRT-based verifiable weighted
        scheme \cite{vwss2025} с весами по правилу $\omega_i = 1/|H_{j(i)}|$.
  \item Выбрать группы для FedGT-style group testing \cite{fedgt2023} совпадающими
        с домохозяйствами — получить byzantine-детекцию \guillemotleft
        бесплатно\guillemotright{} поверх уже существующей группировки.
  \item Замерить фактический размер доли на робота против границы
        $\frac{|S|+|K|}{t}$ \cite{optcss2025} как критерий эффективности реализации.
\end{enumerate}

\textbf{Статус.} Это по-прежнему предлагаемый план, не проверенная конструкция —
но, в отличие от раздела выше, теперь построен из компонентов 2023--2025 годов,
каждый из которых уже формально доказан и эмпирически оценён по отдельности; здесь
предлагается их совмещение специально под флот роботов, сгруппированных по
домохозяйствам, чего в найденной литературе нет.

\section{Тезаурус}

\begin{description}
  \item[PRSS (Pseudorandom Secret Sharing)] Метод не-интерактивной (без
    коммуникации на каждом раунде) генерации свежих Shamir-разделений
    псевдослучайного значения (в частности, нуля) из заранее распределённых
    PRF-ключей.
  \item[PRF (Pseudo-Random Function)] Функция, генерирующая псевдослучайный выход
    по ключу; здесь используется для детерминированной, но непредсказуемой
    генерации масок на каждый раунд.
  \item[Sum-to-zero (share of zero)] Свойство набора долей: их сумма по всем
    участникам тождественно равна нулю — используется для маскирования без
    искажения итоговой суммы полезных значений.
  \item[Straggler-resilient] Устойчивый к тому, что часть участников не успевает
    /не выходит на связь в рамках раунда протокола.
  \item[Threshold Paillier] Вариант аддитивно гомоморфного шифрования Paillier, в
    котором секретный ключ расшифровки разделён по Shamir среди участников — для
    расшифровки нужна коалиция не менее $t$ из них.
  \item[Аддитивный гомоморфизм] Свойство шифра: $\mathrm{Enc}(m_1)\cdot
    \mathrm{Enc}(m_2) = \mathrm{Enc}(m_1+m_2)$ — позволяет складывать значения, не
    расшифровывая их по отдельности.
  \item[Compartmented (multipartite) access structure] Структура доступа
    секрет-шеринга, в которой участники разбиты на непересекающиеся группы
    (compartments) со своими внутренними порогами.
  \item[Домохозяйственный след $h(S)$] Число различных домохозяйств, представленных
    в подмножестве участников $S$ — предлагаемая единица измерения реальной силы
    коалиции вместо сырого числа роботов $|S|$.
  \item[Weighted Secret Sharing (WSS)] Обобщение порогового секрет-шеринга, где
    участникам присваиваются веса $\omega_i$, и безопасность держится, пока
    суммарный вес скомпрометированных участников меньше порога.
  \item[Verifiable Secret Sharing (VSS)] Секрет-шеринг с доказуемой корректностью
    долей — защищает и от злонамеренных участников, и от злонамеренного дилера, не
    только от honest-but-curious.
  \item[Group testing] Метод идентификации \guillemotleft плохих\guillemotright{}
    элементов небольшой группой пересекающихся тестов вместо проверки каждого
    элемента по отдельности; в FedGT — способ найти вредоносных клиентов по
    агрегатам пересекающихся групп.
  \item[Key-homomorphic masking] Совмещение маскирования и гомоморфного шифрования,
    при котором маска шифруется вместе с полезным значением, а не заменяет
    шифрование или заменяется им.
  \item[Non-colluding servers] Модель угроз с несколькими серверами, в которой
    предполагается, что они не сговариваются между собой; безопасность держится,
    пока хотя бы один сервер честен (Prio) или пока не превышен порог сговора
    (ABY3-стиль честного большинства).
  \item[Replicated secret sharing] Схема, в которой каждому из $m$ серверов
    достаётся не одна, а несколько (например, 2 из 3) долей значения — даёт
    устойчивость к выходу части серверов из строя.
\end{description}

\begin{thebibliography}{9}
\bibitem{prss2005} R.~Cramer, I.~Damgård, Y.~Ishai. \emph{Share Conversion,
  Pseudorandom Secret-Sharing and Applications to Secure Computation}. TCC, LNCS
  3378, pp.~342--362, 2005.
\bibitem{prssstraggler2021} F.~Benhamouda, E.~Boyle, N.~Gilboa, S.~Halevi,
  Y.~Ishai, A.~Nof. \emph{Generalized Pseudorandom Secret Sharing and Efficient
  Straggler-Resilient Secure Computation}. TCC 2021; IACR ePrint 2021/1223.
\bibitem{bonawitzref} K.~Bonawitz et al. \emph{Practical Secure Aggregation for
  Privacy-Preserving Machine Learning}. ACM CCS, 2017.
\bibitem{paillier1999} P.~Paillier. \emph{Public-Key Cryptosystems Based on
  Composite Degree Residuosity Classes}. EUROCRYPT, 1999.
\bibitem{thresholdpaillier} P.-A.~Fouque, G.~Poupard, J.~Stern. \emph{Sharing
  Decryption in the Context of Voting or Lotteries}. FC 2000, LNCS 1962,
  pp.~90--104, 2001; I.~Damgård, M.~Jurik. \emph{A Generalisation, a
  Simplification and Some Applications of Paillier's Probabilistic Public-Key
  System}. PKC 2001, LNCS 1992, pp.~119--136, 2001.
\bibitem{compartmented1990} G.~Simmons. \emph{How to (Really) Share a Secret}.
  CRYPTO, 1990; E.~Brickell. \emph{Some Ideal Secret Sharing Schemes}. EUROCRYPT,
  1989/1990; H.~Ghodosi, J.~Pieprzyk, R.~Safavi-Naini. \emph{Secret Sharing in
  Multilevel and Compartmented Groups}. ACISP 1998, LNCS 1438, Springer, 1998.
\bibitem{vwss2025} K.~Shehata, F.~Han, S.A.~Thyagarajan. \emph{Verifiable Weighted
  Secret Sharing}. arXiv:2505.24289, 2025.
\bibitem{fedgt2023} M.~Xhemrishi, J.~Östman, A.~Wachter-Zeh, A.~Graell i Amat.
  \emph{FedGT: Identification of Malicious Clients in Federated Learning with
  Secure Aggregation}. arXiv:2305.05506, 2023.
\bibitem{optcss2025} I.L.~Aureliano, A.~Cohen, R.G.L.~D'Oliveira. \emph{Optimal
  Computational Secret Sharing}. arXiv:2502.02774, 2025.
\bibitem{lerna2023} H.~Li, H.~Lin, A.~Polychroniadou, S.~Tessaro. \emph{LERNA:
  Secure Single-Server Aggregation via Key-Homomorphic Masking}. ASIACRYPT 2023.
\bibitem{dmsafl2026} \emph{DMSA-FL: Secure and Robust Federated Learning via
  Double-Masked Secure Aggregation with Multiparty Homomorphic Encryption}.
  Journal of King Saud University -- Computer and Information Sciences, 2026.
\bibitem{vmhesa2025} \emph{V-MHESA: A Verifiable Masking and Homomorphic
  Encryption-Combined Secure Aggregation Strategy for Privacy-Preserving
  Federated Learning}. Mathematics, DOI:10.3390/math13223687, 2025.
\bibitem{prio2017} H.~Corrigan-Gibbs, D.~Boneh. \emph{Prio: Private, Robust, and
  Scalable Computation of Aggregate Statistics}. USENIX NSDI, 2017.
\bibitem{aby3_2018} P.~Mohassel, P.~Rindal. \emph{ABY3: A Mixed Protocol
  Framework for Machine Learning}. ACM CCS, 2018.
\end{thebibliography}

\end{document}
